% ============================================================================
% PolicyWatcher: arXiv submission source (expanded)
% Compile: pdflatex (2 passes) or tectonic. Standard TeX Live packages only,
% so the file compiles unmodified on arXiv and Overleaf.
% Suggested arXiv categories: cs.CY (primary), cs.SE and cs.CR (cross-list).
% Style note: this document deliberately avoids em dashes and en dashes.
% ============================================================================
\documentclass[11pt]{article}

\usepackage[utf8]{inputenc}
\usepackage[T1]{fontenc}
\usepackage[margin=1in]{geometry}
\usepackage{microtype}
\usepackage{graphicx}
\usepackage{booktabs}
\usepackage{tabularx}
\usepackage{enumitem}
\usepackage{amsmath}
\usepackage{amssymb}
\usepackage{amsthm}
\usepackage{url}
\usepackage[hidelinks]{hyperref}
\usepackage{xcolor}
\usepackage{listings}
\usepackage{placeins}

\setlength{\emergencystretch}{2em}

\lstset{
  basicstyle=\ttfamily\small,
  breaklines=true,
  frame=single,
  framesep=4pt,
  columns=fullflexible,
}

\newcommand{\pw}{\textsc{PolicyWatcher}}
\newtheorem{invariant}{Invariant}

\title{\pw: An Evidence-Preserving Pipeline for Monitoring\\
Terms-of-Service and Privacy-Policy Changes\\
on Consumer Technology Platforms}

\author{
  Fabrizio Degni\\
  \small PolicyWatcher Project\\
  \small \url{https://policywatcher.online}
}

\date{30 July 2026}

\begin{document}
\maketitle

\begin{abstract}
Terms of service and privacy policies govern the digital lives of billions of
people, yet they change silently. Providers routinely amend clauses on data
collection, artificial-intelligence training, and content licensing without
effective notice to the users bound by them. We present \pw, an open, low-cost
platform that acquires configured policy pages, detects text changes, and
publishes bilingual plain-language analyses only after an explicit evidence
gate. The central contribution is an \emph{evidence-preserving monitoring
architecture} that separates source acquisition, immutable snapshots,
AI-assisted interpretation, publication eligibility, and downstream reuse. The
current system combines (i) a five-strategy retrieval cascade spanning direct
HTTP, explicit HTTP/2, remote headless rendering, and two freshness-guarded web
archives; (ii) content and transport validation for bot walls, soft not-found
pages, consent interstitials, redirect drift, and network egress risk; (iii)
scan-level acquisition deduplication with persisted run, retrieval, remediation,
and historical-reference records; and (iv) a first-baseline protocol that
promotes exact source-supported content without emitting a false change, model
score, or notification. Evidence-governance packets, deterministic collection
exports, a forward-only public change feed, and a local webhook receiver
conformance lab extend the same provenance boundary into review and integration
workflows. A dated public snapshot collected on 30 July 2026 reports 50 policy
records, of which 45 passed the public evidence gate and five were withheld;
the public change inventory contained 35 records. These figures describe one
operational state, not semantic accuracy, legal compliance, market coverage, or
adoption. The release was validated by 485 automated tests across 89 test files,
a production build, static checks, a read-only source inventory, and a dry-run
baseline-repair control. The platform runs on commodity shared hosting plus a
small virtual private server and is released as inspectable source code. We
discuss the threat model, public-interest retrieval ethics, AI transparency,
limits of the current evidence claims, and a protocol for longitudinal and
expert-annotated evaluation.
\end{abstract}

\section{Introduction}
\label{sec:intro}

The notice-and-consent model remains the legal backbone of the relationship
between consumers and online platforms, despite two decades of evidence that it
does not function as notice for ordinary users. Reading the privacy policies a
person encounters in a single year has been estimated to require on the order of
two hundred hours \cite{mcdonald2008cost}. Users overwhelmingly accept terms
unread, including clauses that no reasonable person would accept if they were
read \cite{obar2020biggest}, and the language of the documents themselves is
frequently pitched above the reading level of the general public
\cite{fabian2017readability, reidenberg2015disagreeable}. The problem is
compounded by change. Policies are living documents, and the clauses that matter
most to citizens, including data retention, cross-border transfer,
artificial-intelligence training on user content, and licensing of
user-generated works, are precisely the clauses that providers have amended most
aggressively during the generative-AI transition. A user who consented in
January is bound in June by text they have never seen, and the mechanisms that
would surface such a change to the affected population do not exist at scale.

Civil-society initiatives have partly addressed the static version of this
problem. ToS;DR crowdsources human-authored summaries and ratings of terms of
service \cite{tosdr}, and the Electronic Frontier Foundation's TOSBack project
pioneered automated tracking of terms pages \cite{tosback}. On the research side,
a substantial literature has produced annotated corpora \cite{wilson2016creation},
million-document longitudinal archives \cite{amos2021privacy}, and automated
clause-level analysis using natural-language processing and machine learning
\cite{harkous2018polisis, srinath2021privacy, lippi2019claudette,
andow2019policylint, cui2023poligraph}. Yet a gap persists between
\emph{corpus construction}, which is retrospective and research-oriented, and
\emph{civic monitoring}, which must be continuous, forward-looking, and
communicated in plain language to the people affected.

Closing that gap requires solving an engineering problem that is easy to
underestimate. \textbf{Reliably obtaining the current, complete text of a policy
page is adversarial.} The pages that matter most sit behind bot-mitigation
walls, JavaScript-only rendering, geographic content splits, and frequent URL
churn. Naive scrapers silently capture error pages, consent interstitials, or
stale cached copies, and every downstream artifact, from a stored diff to a risk
score to a subscriber email, inherits that corruption without any signal that it
has occurred. The dominant failure mode in this domain is not a crash. It is a
plausible wrong value that propagates as if it were fact.

\pw{} is a running platform that monitors 50 policy documents across 16
technology and fintech providers in up to four jurisdictional variants
(European Union, United States, United Kingdom, and Global), detects changes,
and publishes bilingual (English and Italian) analyses, risk assessments,
comparison matrices, and subscriber notifications. This paper reports its design
as a systems contribution, organized around a single question:

\begin{quote}
\emph{How can a public-interest monitoring system keep each public record tied
to retrieved evidence, and withhold the record when that condition is not met?}
\end{quote}

\noindent The contributions of this paper are the following.

\begin{enumerate}[leftmargin=*]
  \item \textbf{An evidence-grade acquisition cascade} (Section~\ref{sec:pipeline})
  that combines direct fetching with realistic browser fingerprints, explicit
  HTTP/2 negotiation for protocol-strict providers, remote headless rendering
  for script-dependent pages, and two web archives, all governed by a
  \emph{freshness guard} with a stated timestamp invariant that rejects archive
  snapshots captured before the latest live baseline (Section~\ref{sec:freshness}).
  \item \textbf{A two-layer validation model} (Section~\ref{sec:validation})
  that separates transport-level outcomes from content-level pathologies,
  including bot challenges, soft not-found pages, maintenance pages, paywalls,
  consent walls, and redirect drift, behind a deterministic result contract that
  forces every caller to handle failure explicitly, and an egress-side
  server-side request forgery guard with DNS pinning.
  \item \textbf{A source-reliability control plane}
  (Section~\ref{sec:reliability}) that normalizes acquisition keys, deduplicates
  shared retrievals, persists scan and retrieval counts, records structured
  failure causes, and retains stale archive candidates only as historical
  references excluded from change detection.
  \item \textbf{A data-governance layer} (Section~\ref{sec:governance}) that
  encodes the distinction between fixtures and evidence in the database schema,
  gates every public surface on evidence status, and defines an exact-evidence
  first-baseline protocol that emits no spurious change event, model score, or
  notification.
  \item \textbf{Citable evidence and bounded integration contracts}
  (Section~\ref{sec:delivery}) comprising change-bound evidence packets,
  deterministic browser-local collections and handoff manifests, a forward-only
  public change-event feed, and a receiver conformance suite whose results are
  explicitly separated from endpoint certification or delivery guarantees.
  \item \textbf{A threat model and an adversarial self-assessment}
  (Sections~\ref{sec:threat} and~\ref{sec:redteam}) conducted against the
  running system, framed by two questions: where can uncertain data still become
  public evidence, and which function can destroy or promote data irreversibly
  if an operational state is wrong. We report the findings and the guardrails
  that closed them.
  \item \textbf{An implementation and reproducibility account}
  (Section~\ref{sec:impl}) covering the open-source release, the automated test
  suite, continuous integration, and third-party quality attestation, together
  with a deployment blueprint (Section~\ref{sec:deployment}) that runs the entire
  platform on commodity shared hosting plus one small virtual private server.
\end{enumerate}

\section{Background and Related Work}
\label{sec:related}

\subsection{The readability and change problem}
The mismatch between the legal function of policies and their practical
readability is well documented. McDonald and Cranor quantified the aggregate
time cost of reading policies \cite{mcdonald2008cost}; Reidenberg et al.\
demonstrated systematic mismatches between policy meaning and user understanding
\cite{reidenberg2015disagreeable}; and large-scale readability analyses confirm
that policies routinely exceed the comprehension level of the general population
\cite{fabian2017readability}. Obar and Oeldorf-Hirsch showed experimentally that
users accept unread terms at near-universal rates \cite{obar2020biggest}. These
results motivate a monitoring system oriented toward the affected public rather
than toward researchers, and they motivate a focus on \emph{change}, because a
reader who cannot be expected to read a policy once certainly cannot be expected
to re-read it after every amendment.

\subsection{Policy corpora and longitudinal analysis}
Wilson et al.\ introduced OPP-115, the first richly annotated privacy-policy
corpus, establishing a data-practice taxonomy that later systems reused
\cite{wilson2016creation}. Amos et al.\ curated more than one million policy
snapshots spanning two decades by mining the Wayback Machine, and documented both
the growth of policy length over time and the methodological hazards of
archive-based curation, notably error pages and boilerplate masquerading as
policies \cite{amos2021privacy}. Linden et al.\ analyzed the policy landscape
before and after the General Data Protection Regulation and found measurable
changes in length and disclosure \cite{linden2020landscape}. The content
validation layer in \pw{} (Section~\ref{sec:validation}) operationalizes the
hazards these corpora identified as run-time checks, and the freshness guard
(Section~\ref{sec:freshness}) addresses a failure mode specific to \emph{live}
monitoring that retrospective corpora, which fix a snapshot set in advance, do
not face.

\subsection{Automated policy and terms analysis}
A rich line of work applies natural-language processing to policies and terms.
Privee and Polisis demonstrated automated classification and question answering
over privacy policies \cite{zimmeck2014privee, harkous2018polisis}; PrivaSeer
scaled corpus construction to the open web \cite{srinath2021privacy}; PolicyLint
and PoliGraph extracted and reasoned over data-flow statements, including
internal contradictions \cite{andow2019policylint, cui2023poligraph}; and MAPS
and related systems scaled compliance analysis to mobile-app ecosystems
\cite{zimmeck2019maps}. In the terms-of-service setting specifically, the
CLAUDETTE line of work trained detectors of potentially unfair contractual
clauses \cite{lippi2019claudette}. \pw{} consumes a general-purpose large
language model for bilingual summarization and risk annotation, but treats the
model as an \emph{untrusted component}: its outputs are schema-validated,
range-checked, and, by policy, never a substitute for retrieval evidence
(Section~\ref{sec:analysis}). The system's contribution is therefore
complementary to this literature. It is not a new clause classifier but an
infrastructure that supplies the classifier, whatever it is, with a validated
retrieval path and publishes its outputs only when the provenance condition is
met.

\subsection{Consent interfaces and web measurement}
A parallel measurement literature studies the interfaces that gate access to
policy-governed services. Studies of cookie consent notices and their dark
patterns \cite{utz2019consent, nouwens2020dark, matte2020cookie,
sanchezrola2019optout} document precisely the consent walls and interstitials
that a policy scraper must detect and reject rather than ingest. Work on
realistic and reproducible web crawling \cite{jueckstock2021realistic} and on
the detection of browser fingerprinting \cite{iqbal2021fingerprinting} frames the
technical arms race that makes reliable retrieval adversarial and that motivates
the multi-strategy cascade in Section~\ref{sec:pipeline}. Automated auditing of
third-party disclosure in policies \cite{libert2018auditing} further underscores
that a policy's text and a site's actual behavior can diverge, which is one
reason \pw{} anchors every published record to a specific retrieval event.

\subsection{Monitoring services and positioning}
TOSBack \cite{tosback} and ToS;DR \cite{tosdr} are the closest antecedents in
spirit. TOSBack tracks raw page text and TOSBack-style diffs; ToS;DR provides
human-curated ratings. \pw{} differs in three respects. First, it invests in
adversarial retrieval, which is required because modern policy pages actively
resist naive fetching in ways that were less prevalent when earlier trackers
were designed. Second, it maintains a per-document provenance model in which
every published datum links to a check log that records the transport used, the
HTTP status, the content hash, and the source of retrieval. Third, it governs
the boundary between demonstration data and evidence explicitly, which we found
in practice to be the dominant integrity risk (Section~\ref{sec:redteam}).
Regulatory context is provided by the transparency obligations of the General
Data Protection Regulation \cite{gdpr} and the terms-and-conditions provisions
of the Digital Services Act \cite{dsa}; independent monitoring gives such
obligations empirical traction.

The EU Artificial Intelligence Act introduces transparency obligations for
providers and deployers of certain AI systems in Article 50
\cite{aiact}. The European Commission's July 2026 guidelines state that these
obligations apply from 2 August 2026 and provide practical guidance on scope and
implementation \cite{ecarticle50}. \pw{} does not assess conformity with those
obligations. The relevance is architectural: AI-assisted public-interest text,
source provenance, human-review boundaries, and machine-readable distribution
should be distinguishable and inspectable rather than presented as a single
authoritative output.

\section{System Overview}
\label{sec:overview}

\pw{} is implemented as a server-rendered web application (Next.js with the App
Router, React, and TypeScript) backed by a relational database accessed through
an object-relational mapper (Prisma over SQLite). Policy analysis calls a
hosted large language model. Headless rendering, which requires a Chromium
runtime that shared hosting cannot provide, runs as a separate authenticated
service (Section~\ref{sec:deployment}). The platform comprises five subsystems.

\begin{enumerate}[leftmargin=*]
  \item \textbf{Acquisition}: a scheduled scanner that, for each configured
  policy URL, executes the retrieval cascade and validation of
  Sections~\ref{sec:pipeline} and~\ref{sec:validation}, producing either a
  normalized text with a SHA-256 content hash or an explicit failure status.
  \item \textbf{Change detection and analysis}: hash comparison against the
  stored baseline, and, on change, a stored full-text line diff plus a
  model-generated analysis comprising bilingual summaries, key points,
  remediation guidance, a risk score, a fifteen-indicator assessment, and
  per-region impact statements (Section~\ref{sec:analysis}).
  \item \textbf{Publication}: a public dashboard, a change timeline, per-change
  pages with diffs, a cross-company comparison matrix over the indicator
  taxonomy, regional heatmaps and radar benchmarks, an evidence-ranked signals
  board, a regulatory observatory, embeddable widgets, downloadable reports, and
  subscriber email digests, all gated on evidence status
  (Section~\ref{sec:governance}).
  \item \textbf{Evidence reuse and integration}: change-bound evidence packets,
  deterministic evidence collections, vendor-neutral handoff manifests, a
  public change-event polling feed, versioned JSON Schemas, and a local receiver
  verification and conformance workbench (Section~\ref{sec:delivery}).
  \item \textbf{Operations and audit}: an administrative surface exposing
  per-policy check logs, scan runs, source retrievals, remediation state,
  dataset-quality findings with severity-ranked blockers, explainability review,
  an append-only review log, an access log, and health and control endpoints for
  the remote renderer.
\end{enumerate}

\paragraph{Data model.}
Persistence is separated into twenty relational entities. \texttt{Company} and \texttt{Policy} hold
the configured monitoring scope; a policy carries its source URL, jurisdiction,
current text and hash, a lifecycle value, and the ingestion method that records
how the current text was obtained. A policy snapshot stores versioned full texts,
each flagged as public evidence. A policy-change record stores the diff and the
analysis fields. A policy-check log stores one row per retrieval attempt,
recording status, source, HTTP status, final URL, content hash, text length,
failure reason, and, for archive retrievals, the snapshot timestamp.
\texttt{ScanRun}, \texttt{SourceRetrieval}, \texttt{SourceRemediationIssue}, and
\texttt{HistoricalSourceReference} retain scan-level accounting, shared
acquisition state, repeated failures, and dated archive context without merging
them into public evidence. Additional entities cover discovery and onboarding,
policy inquiries, dataset-QA decisions, administrative review, access logging,
press measurement, regional impact, and subscribers. The check log and source
retrieval records form the provenance spine: public records retain a link to the
attempt that produced their source evidence while operational telemetry remains
separately queryable.

\section{The Evidence Acquisition Cascade}
\label{sec:pipeline}

\subsection{Design invariant: evidence or explicit status}
The pipeline's contract is that a retrieval attempt returns exactly one of three
statuses. \texttt{ok} carries validated text and a hash. \texttt{unavailable}
signals a transient failure such as a timeout, a bot block, or a maintenance
page. \texttt{invalid} signals that the URL no longer denotes the intended
document, for example a hard not-found response, a soft not-found page, or a
redirect that drifts to an unrelated destination. Nothing else is
representable in the result type. A caller cannot receive partial text without a
hash, or a hash without provenance. Failure paths write a check-log row and a
user-facing notice that directs the reader to the official source, rather than
an approximation. This shape was chosen after observing that most data-integrity
incidents in scraping systems are not exceptions but plausible wrong values that
propagate silently.

\subsection{The retrieval cascade}
Strategies are ordered by fidelity to the live site. Archives come last because
they can only witness the past. Table~\ref{tab:cascade} summarizes the cascade.

\begin{table}[htbp]
\centering
\small
\begin{tabularx}{\textwidth}{@{}l l X@{}}
\toprule
\textbf{Order} & \textbf{Strategy} & \textbf{Role and rationale} \\
\midrule
1 & Direct HTTP/1.1 & Realistic desktop browser fingerprint; manual redirects with per-hop validation; body-aware timeout. Primary path. \\
2 & Explicit HTTP/2 & Direct ALPN negotiation for providers that reject HTTP/1.1 from non-browser clients with protocol-level errors. \\
3 & Rendered fetch & Remote headless Chromium executes JavaScript; recovers single-page-application shells and many bot-protected pages. Optional. \\
4 & Wayback Machine & Freshness-guarded archived snapshot; witness of last resort. \\
5 & Common Crawl & Freshness-guarded archive; gzip-compressed records are decompressed before extraction. \\
\bottomrule
\end{tabularx}
\caption{The five-strategy retrieval cascade. Strategies 1 to 3 observe the live
site; strategies 4 and 5 are archives subject to the freshness guard of
Section~\ref{sec:freshness}.}
\label{tab:cascade}
\end{table}
\FloatBarrier

\textbf{Direct fetch.} The primary path issues an HTTP/1.1 request with a
realistic desktop browser fingerprint (User-Agent, client-hint headers, and
fetch-metadata headers), follows redirects manually with server-side request
forgery re-validation at every hop, and retries on status 429 and 5xx with
exponential backoff. The request timeout of twenty seconds covers the response
\emph{body}, not merely the response headers, which matters against slow-drip
responses that deliver headers promptly and then stall. Responses whose content
type is clearly not HTML, such as PDF, images, or octet streams, are rejected
rather than parsed. The transport makes up to three attempts and rotates among
three browser profiles.

\textbf{Explicit HTTP/2.} Several large providers reject HTTP/1.1 requests from
non-browser clients with protocol-level 400 responses. A dedicated client
negotiates HTTP/2 over ALPN directly. Response chunks are accumulated as bytes
and decoded once, because decoding at arbitrary chunk boundaries corrupts
multi-byte UTF-8 sequences, and redirects are treated as failures at this layer
by construction.

\textbf{Rendered fetch.} Pages that ship as JavaScript application shells, which
is increasingly the norm for exactly the providers most relevant to consumers,
are delegated to a self-hosted rendering service (Section~\ref{sec:deployment}).
The service loads the page in headless Chromium, waits for hydration and for
network quiescence, and returns the rendered document. If the service is not
configured or is unreachable, the cascade continues to the archives. Rendering
is an amplifier, not a dependency.

\textbf{Freshness-guarded archives.} The Wayback Machine \cite{wayback} and
Common Crawl \cite{commoncrawl} serve as witnesses of last resort. Common Crawl
records are located by CDX lookup and fetched by byte-range reads of
gzip-compressed WARC records, which are decompressed before parsing. Both archive
strategies enforce the freshness guard of Section~\ref{sec:freshness}.

\subsection{Content validation}
\label{sec:validation}
The second layer validates \emph{what} was retrieved. Its checks are pure
functions of the retrieved bytes and are unit-tested in isolation.

\begin{itemize}[leftmargin=*]
  \item \textbf{Block-page detection}: signatures of CAPTCHA and JavaScript
  challenges, web application firewall denial pages, maintenance pages,
  paywalls and login walls, and consent interstitials, matched only against the
  document head so that a policy which legitimately \emph{discusses} CAPTCHAs is
  not itself flagged.
  \item \textbf{Soft not-found detection}: short pages using not-found
  phrasings in multiple languages, which catches servers that return status 200
  for missing documents.
  \item \textbf{Structural extraction}: removal of boilerplate (navigation,
  headers, footers, cookie banners, hidden elements) and conversion of the main
  content region to deterministic, structured plain text. The
  \emph{same} normalizer feeds storage, hashing, and diffing, so cosmetic markup
  changes do not register as policy changes. When a configured URL carries a
  fragment identifier, extraction is scoped to the addressed section.
  \item \textbf{Plausibility checks}: a minimum length of 800 characters and a
  policy-marker vocabulary test that requires at least three of twenty
  domain terms, which rejects well-formed pages that are not policies. Extracted
  text is capped at 500{,}000 characters to bound storage.
  \item \textbf{Drift rejection}: a live retrieval whose final URL, after
  redirects, lands on a different registrable domain (host drift) or on a site
  root when a document path was configured (path drift) is rejected as
  \texttt{invalid} rather than ingested. This closes the failure path in which a
  redirect to a homepage whose footer contains the word ``privacy'' would
  otherwise pass the marker test.
\end{itemize}

\subsection{Egress protection}
Because the scanner fetches administrator-configured URLs from server-side
infrastructure, every fetch passes a server-side request forgery guard \cite{owasp2021}: a scheme
allowlist, rejection of credentials embedded in the URL, DNS resolution with
blocking of private, loopback, and link-local addresses, and re-validation on
every redirect hop. To close the time-of-check-to-time-of-use gap between
validation and connection, the HTTP/1.1 and HTTP/2 clients pin the validated
address at connection time rather than re-resolving the hostname, and the IPv6
address parser is hardened against compressed and mixed notations. The companion
renderer applies the same address policy to every browser-initiated subresource
request and checks redirect coherence against the Public Suffix List.

\section{The Freshness Guard}
\label{sec:freshness}

The archive strategies introduce a subtle integrity risk that is specific to
live monitoring. Suppose the live site is temporarily blocked. A naive cascade
would fall through to an archive and recover an older copy of the document, whose
hash differs from the stored baseline. The system would then register a
\emph{change}, run analysis, and notify subscribers about a regression to a
superseded text. When the live site later recovers, a second spurious change
fires. We eliminate this oscillation class with a single rule.

\begin{quote}
An archived snapshot is admissible only if its capture timestamp is not earlier
than the policy's last successful live retrieval.
\end{quote}

\noindent Let $t_\ell$ be the timestamp of the last retrieval that produced a
stored baseline, and let $t_s$ be the capture timestamp of a candidate archive
snapshot. The guard admits the snapshot only if $t_s \ge t_\ell$.

\begin{invariant}
Provided that the recorded archive capture timestamp is accurate, the freshness
guard rejects any candidate snapshot captured before the current baseline.
\end{invariant}

\begin{proof}[Proof sketch]
The stored baseline was written by a retrieval at time $t_\ell$, and by the
pipeline invariant it reflects the document observed at $t_\ell$. An archive
snapshot is admitted only when $t_s \ge t_\ell$, so any admitted snapshot was
captured no earlier than the current baseline. A change is recorded only when an
admitted retrieval yields a hash different from the baseline hash. Therefore a
recorded change from an archive source refers to a capture at or after $t_\ell$,
not before it. The rule prevents a pre-$t_\ell$ capture from becoming a new
baseline under those assumptions. Snapshots lacking a timestamp are treated as
inadmissible, so the rule does not rely on an unknown capture time.
\end{proof}

The guard is enforced in both archive strategies. Admitted archive retrievals
additionally carry their capture timestamp into the check log, so that the age of
recovered evidence is recorded and can be displayed to readers. A practical side
effect is a reduction in load on public-interest archive infrastructure, since
stale snapshots are rejected before they are fetched in full.

\section{Source Reliability Control Plane}
\label{sec:reliability}

The acquisition cascade originally persisted policy-level check logs but did
not make the relationship between selected records, unique source acquisitions,
network attempts, and reused results explicit. This matters when several policy
records share a legal hub, or when a canonical policy URL points to a separate
official retrieval mirror or PDF. Treating every row as a network fetch inflates
operational counts and can apply inconsistent archive-freshness floors.

The current scanner computes a normalized \emph{retrieval key} for each selected
policy. Records with the same key share one acquisition while retaining
independent comparisons, snapshots, jurisdictions, and publication decisions.
The strictest last-successful-live timestamp among the sharing records becomes
the archive-freshness floor. A persisted \texttt{ScanRun} records selected rows,
unique retrieval keys, network attempts, deduplicated acquisitions, unavailable
sources, and dependency failures. Each \texttt{SourceRetrieval} records the
normalized key, strategy, outcome, structured cause code, HTTP metadata, hash,
and affected-record count.

Failures are classified by stable causes such as source unavailability, content
rejection, dependency failure, and invalid configuration. Repeated failures can
open a \texttt{SourceRemediationIssue}; the record suggests review actions but
never authorizes bypassing a provider access control. Archive candidates that
fail the freshness guard are stored only as
\texttt{HistoricalSourceReference} records. They can explain source continuity
but cannot create a baseline or a change.

The protected Source Reliability console exposes these distinctions and the
public-baseline coverage count. It is intentionally an operational control
plane, not a public quality score. A zero or partial coverage state means that
the evidence gate is withholding records; it is not evidence of provider
quality, system security, or legal compliance.

\section{Change Detection and Analysis}
\label{sec:analysis}

Change detection is hash equality over normalized text. On inequality, the system
stores a new immutable snapshot and a full-text line diff. An earlier design
truncated the diff to the first ten thousand characters, which silently hid
changes beyond the opening pages of long documents; the current design stores the
complete diff. The system then generates an analysis with a hosted large language
model. The primary model is \texttt{gemini-2.5-flash}; on transient errors such
as rate limiting or temporary unavailability, the system falls back to
\texttt{gemini-3.5-flash-lite}. Generation uses a low temperature, constrains
the response with provider-side JSON Schema, and validates the payload again
locally before normalization or persistence.

The analysis schema contains bilingual executive summaries, a one-sentence
take-away in each language, three to five key points with sentiment, exactly
three risk rationales, an overall risk band, an integer risk score from 1 to 10,
two to four remediation actions with links to official opt-out or objection
mechanisms, exactly six per-region impact statements (the cross product of three
regions and two perspectives), and the fifteen-indicator assessment of
Table~\ref{tab:kpi}.

\begin{table}[htbp]
\centering
\small
\begin{tabularx}{\textwidth}{@{}l X@{}}
\toprule
\textbf{Group (5 indicators each)} & \textbf{Indicators and representative values} \\
\midrule
Privacy and data protection &
Data collection (Minimal / Moderate / Extensive); Third-party sharing
(Restricted / Limited / Broad); Data retention (Defined / Extended /
Indefinite); Right to deletion (Full / Partial / Not available); Cross-border
transfer (Restricted / Controlled / Unrestricted). \\
\addlinespace
AI governance &
Training opt-out (Available / Opt-out / Not available); Output ownership
(User retained / Shared / Company retained); Algorithmic transparency
(Published / Mentioned / Opaque); Automated decisions (Transparent / Partial /
Opaque); Bias and fairness (Committed / Mentioned / Absent). \\
\addlinespace
Ethics and corporate governance &
Consent mechanism (Explicit opt-in / Opt-out / Implicit); Regulatory compliance
(Comprehensive / Partial / Minimal); Breach notification (Within 24h / Within
72h / Unspecified); Independent audit (Certified / Mentioned / Absent); Content
moderation (Transparent / Partial / Opaque). \\
\bottomrule
\end{tabularx}
\caption{The fifteen-indicator taxonomy used for cross-company comparison. Every
indicator additionally admits the value ``Not assessed'', which is the default
until an assessment is made.}
\label{tab:kpi}
\end{table}
\FloatBarrier

Three properties keep the model subordinate to evidence. First, analyses are
generated only for retrieved text; there is no path from model output to policy
text. Second, model outputs are parsed against a strict schema with range
validation: an out-of-range risk score or an unrecognized risk band causes the
analysis to be rejected rather than clamped, and indicator values are passed
through a case-insensitive whitelist that maps anything unrecognized to ``Not
assessed''. Third, in production the system prefers to block ingestion of an
analysis over substituting a heuristic mock; the demonstration fallback is
available only when an explicit environment flag is set. The conversational
interface that answers user questions over a policy applies an explicit
prompt-injection guard that treats the policy context as untrusted input and
returns an honest ``unavailable'' response rather than fabricating an answer.

The public analysis layer exposes both unified and side-by-side diffs, exact
snapshot hashes, region-aware heatmaps, and benchmark radar views. Normalized
benchmark values are accompanied by original categorical values, missing-value
states, and a statement that the visualizations are comparative review aids,
not compliance or business-performance scores. Where the stored analysis
contains a source passage, publication requires an exact substring match
against the declared snapshot side. Unmatched passages are withheld rather than
presented as citations. Governance references to the EU AI Act, ISO/IEC 42001,
NIST AI RMF, and OECD AI Principles are advisory review mappings, not legal
interpretations or conformity assessments.

\section{Citable Evidence and Bounded Integration}
\label{sec:delivery}

The platform extends the public evidence gate through four reuse contracts.
First, a change-bound \emph{Evidence Packet} exposes sanitized source-confidence
state, snapshot identifiers, SHA-256 fingerprints, verified source passages,
advisory governance references, human-review questions, and a deterministic
content digest as JSON or PDF. Second, browser-local \emph{Evidence Collections}
group up to twelve exact public changes. Shared URLs contain only sorted public
change identifiers; local titles and review state do not leave the browser. A
collection can be exported as deterministic JSON, Markdown, CSV, or a
vendor-neutral handoff manifest containing evidence links, digests, review
questions, acceptance criteria, and receiving-system instructions.

Third, a versioned public change-event feed provides bounded forward polling of
already-published events. It uses deterministic event identifiers, an opaque
cursor, separate publication time, localized summaries, and exact evidence
links. It is not a subscriber registry and provides no push delivery, recipient
state, retry, or receipt. Fourth, a candidate HMAC-\allowbreak SHA-256 receiver contract is
published with exact headers, signing input, a public test vector, and Node and
Python examples. A browser-local Receiver Conformance Lab executes eight
positive and negative fixtures and exports only case identifiers, expected and
actual decision codes, totals, and the interpretation boundary.

These contracts are deliberately separable. A deterministic digest establishes
artifact identity, not semantic truth. A passing receiver fixture establishes
compatibility with that case, not endpoint identity, production secret custody,
replay protection, delivery reliability, or implementation certification.
Direct publishing to Jira, Confluence, Microsoft Teams, Slack, or a GRC system
remains outside the public contract until identity, authorization, delivery,
retry, and audit controls are implemented.

\section{Functional Surface at Beta 21}
\label{sec:features}

Table~\ref{tab:features} records the main available surfaces so that the systems
claims can be separated from the size of the product interface. Availability
means that a route or contract exists in the inspected release; it does not
establish adoption, completeness, continuous uptime, or independent validation.

\begin{table}[htbp]
\centering
\small
\begin{tabularx}{\textwidth}{@{}l X X@{}}
\toprule
\textbf{Domain} & \textbf{Available functions} & \textbf{Boundary} \\
\midrule
Acquisition and monitoring & Five-path retrieval cascade; complete and bounded scans; normalized shared acquisition keys; source continuity; structured remediation. & External source access and structure remain provider-dependent. \\
\addlinespace
Evidence and analysis & Immutable snapshots; full and side-by-side diffs; bilingual summaries; fifteen indicators; source-anchored reasons; change-bound JSON/PDF Evidence Packets. & AI-assisted outputs require review and are not legal findings. \\
\addlinespace
Public exploration & Evidence Console; Timeline; Signals; Observatory; regional heatmap; radar benchmark; Matrix; Atlas; Feature Atlas; Showcase; release-impact Roadmap. & Missing assessments remain ``Not assessed'' and visual normalization is not a compliance score. \\
\addlinespace
Intake and companion access & ``What changed?'' text intake; browser-local bounded MIME parsing; published Chrome extension for selected page text and links. & No mailbox access, attachment processing, or automatic legal classification is claimed. \\
\addlinespace
Collaboration and integration & Browser-local Evidence Collections; deterministic exports; handoff manifests; public v1 and enterprise-oriented v2 read contracts; change-event polling; receiver verification and conformance fixtures. & No persistent team workspace or outbound third-party delivery is available. \\
\addlinespace
Transparency and reuse & Public methodology, Trust and QA pages, Claim Registry, Press Kit, Data Room, release archive, structured data, citable embeds, Pulse story packs, social cards, corrections and coverage registries. & Editorial and aggregate telemetry does not measure truth, reach, endorsement, or conversion. \\
\addlinespace
Protected governance & Grouped admin navigation; Source Reliability; Cron Manager; Dataset QA; explainability; KPI audit; review and access logs; source onboarding; database and VPS readiness. & Protected operations remain role-scoped and optional metrics cannot become an authentication dependency. \\
\bottomrule
\end{tabularx}
\caption{Available PolicyWatcher functional surface in release 3.9.0 Beta 21,
with the principal interpretation boundary for each domain.}
\label{tab:features}
\end{table}
\FloatBarrier

\section{Data Governance: Fixtures versus Evidence}
\label{sec:governance}

A monitoring platform needs demonstration data before its first production scan,
and it needs development fixtures forever after. The dominant integrity risk we
encountered in practice was not a scraping failure but \emph{fixture leakage}:
hand-authored demonstration records, complete with plausible dates and risk
scores, reaching public surfaces indistinguishably from evidence. \pw{} addresses
this structurally.

\paragraph{Schema-level distinction.}
Every policy carries an \texttt{ingestionMethod}, which is either \texttt{Seeded}
or one of the retrieval sources (\texttt{direct}, \texttt{http2},
\texttt{rendered}, \texttt{wayback}, \texttt{commoncrawl}), and a
\texttt{dataStatus} lifecycle value drawn from \texttt{Configured},
\texttt{Available}, \texttt{Partial}, \texttt{Needs Review}, \texttt{Unavailable},
and \texttt{Reviewed}. Snapshots and changes additionally carry a
\texttt{publicEvidence} boolean. Fixtures are born \texttt{Seeded} and
\texttt{Configured}; they are not treated as retrieved evidence.

\paragraph{Public data gate.}
Every public surface, which includes the read APIs, the server-rendered pages,
the share and embed routes, the OpenGraph image routes, the downloadable
reports, the sitemap, and, critically, the subscriber emails, filters on
evidence status. To be public, a policy must have a status in the set
$\{\texttt{Available}, \texttt{Reviewed}\}$, must not be \texttt{Seeded}, and
must have at least one snapshot flagged as public evidence; changes must
themselves be flagged as public evidence. Policies in a suspended state are
represented publicly only by a minimal ``source under verification'' notice, with
no text, no scores, and no timeline events. The gate is a single shared predicate
applied uniformly; Section~\ref{sec:redteam} explains why centralization matters.

A consequence of this design, which we state plainly because it is unusual, is
that the bundled local development database can contain 50 configured policy
rows while admitting none of them to public routes. The records exercise the
interface but remain \texttt{Seeded} and \texttt{Configured} until retrieval
evidence exists. Production and local fixture states must therefore be measured
separately; the cardinality of a configured inventory is never a claim about
public evidence coverage.

\paragraph{First-baseline and repair protocol.}
Promoting fixtures or previously misclassified records to evidence poses a trap.
The first real retrieval necessarily differs from a placeholder, and a naive
pipeline would emit a wave of fabricated ``changes'', each with a model score
and notification. The first-baseline path treats successful source text as a
baseline when no exact verified baseline exists. It stores or promotes the
matching snapshot, updates the policy provenance, and emits no change record,
analysis, or notification.

Existing deployments use a separate dry-run-first repair command. A row is
eligible only when a stored snapshot hash matches a successful source check for
that policy. The dry run lists eligible and retained-private rows; mutation
requires an explicit \texttt{--apply} flag. The repair cannot infer evidence,
promote a stale historical reference, bypass source access controls, or create a
change event. This narrower authorization rule replaces reliance on a mutable
status flag and preserves genuine history.

\paragraph{Quality gates as release blockers.}
A dataset-assurance tool audits the database for source-evidence defects and
treats them as blockers rather than warnings, because the intended production
state is to remain publicly silent until evidence exists. The command-line tool
uses two severities, blocker and warning, and fails the build if any blocker is
present. The administrative console applies a richer three-level scheme (critical,
warning, informational) across roughly twenty issue areas, among them Source Fit,
which flags a global policy whose URL uses a regional path segment or an
EU-or-US policy pointing at a locale-specific page; Coverage, which flags missing
or very short text and missing analysis; Freshness, which flags policies not
successfully checked in more than thirty days; and Hash Integrity, Snapshot
Chain, and Public Evidence Gate checks. Each finding carries a stable identifier
and is reconciled against an append-only review table, so that a human decision to
accept or ignore a finding is itself recorded with its author and rationale.

\section{Threat Model}
\label{sec:threat}

We consider three classes of adversary and one class of operator error. The
first adversary is the \emph{monitored provider} that, deliberately or not,
serves bot walls, geo-splits, or JavaScript-only shells that frustrate
retrieval; the cascade and validation layers are the response. The second is the
\emph{network position} that could induce the server-side fetcher to reach
internal infrastructure; the server-side request forgery guard with DNS pinning
is the response. The third is the \emph{content adversary} that embeds
instructions in policy text intended to steer the language model; the analysis
layer treats retrieved text as untrusted and validates model outputs against a
schema. The fourth class, which our self-assessment found to be the most
consequential in practice, is \emph{operator error}: a wrong status, a
mis-scoped deletion, or a deployment artifact that overwrites live data. The
governance layer, the evidence-based re-baseline predicate, and the operational
guardrails of Section~\ref{sec:redteam} address this class. We explicitly do not
attempt to cover a compromised hosting environment, a leaked administrative
secret used within its validity window, or an archive that itself serves
corrupted data without an observable indication.

\section{Adversarial Self-Assessment}
\label{sec:redteam}

Before the first production re-baseline, we conducted a structured red-team
review of the running system, organized around two questions:

\begin{enumerate}[leftmargin=*]
  \item Where can uncertain, seeded, incomplete, or badly-retrieved data still
  become public evidence, a notification, a timeline entry, a risk score, or
  apparently reliable history?
  \item Which function can delete, overwrite, or promote data irreversibly if an
  operational state is wrong?
\end{enumerate}

The review combined manual code audit, live probing of external dependencies,
database inspection, and parallel automated code exploration. Table~\ref{tab:red}
summarizes representative confirmed findings, all subsequently closed. We report
the methodology and the failure classes rather than dwelling on individual
defects, because the two guiding questions generalize to any civic-evidence
platform and, in our experience, surfaced cross-cutting risks that
component-level checklists did not.

\begin{table}[htbp]
\centering
\small
\begin{tabularx}{\textwidth}{@{}l X X@{}}
\toprule
\textbf{Class} & \textbf{Finding} & \textbf{Guardrail} \\
\midrule
Ungated egress & Periodic email digests queried changes by recency alone,
bypassing the gate that protected web surfaces, so fixture changes could have
been mailed to subscribers as real intelligence. & Apply the same evidence
predicate to every egress channel, including email, feeds, and the sitemap; a
time window is not a filter. \\
\addlinespace
Enumeration & The sitemap enumerated all change identifiers, including gated
ones, advertising non-public resources to crawlers. & Gate the sitemap on the
same predicate as the pages it references. \\
\addlinespace
Archive regression & A transient block on a live site could reinstate an
outdated archived copy as a change; a legacy cache fallback had also become dead
weight upstream. & The freshness guard of Section~\ref{sec:freshness}; removal of
the retired fallback. \\
\addlinespace
Irreversible operation & Re-baseline, cascade deletion of a company, and the
inclusion of a database file in a deployment artifact each offered a single-step
path to irreversible loss under operator error. & Gate destruction on evidence,
not on a flag; exclude data stores from deployment artifacts; separate the
session-signing secret from the API secret; require an encrypted export before
promotion. \\
\bottomrule
\end{tabularx}
\caption{Representative findings from the adversarial self-assessment and the
guardrails that closed them.}
\label{tab:red}
\end{table}
\FloatBarrier

\section{Implementation, Testing, and Reproducibility}
\label{sec:impl}

\pw{} is released as open source. The retrieval and validation logic is factored
so that its decision functions, including URL and address validation, block-page
and soft not-found detection, text extraction, the freshness guard, and the
drift checks, are pure and unit-tested; the public-data gate and the governance
predicates are likewise unit-tested. At the time of writing the suite comprises
485 tests across 89 test files and passes in full. Network-dependent paths, which are
not amenable to deterministic unit testing, are exercised by integration smoke
tests against live sources. Continuous integration runs five workflows on each
change: a quality gate (schema validation, build, lint, dependency audit, and the
dataset-assurance check), a coverage workflow that uploads to a coverage service,
static analysis with CodeQL, a SonarCloud scan, and an OpenSSF Scorecard run.
The release also completed TypeScript validation, linting with no release-source
errors, a production build that generated 130 application routes and pages, a
read-only source-inventory audit, and a dry-run baseline repair. A production
dependency advisory check reported no known advisories at the time of the run;
this is not a claim that the application has no vulnerabilities. The project
additionally holds an OpenSSF Best Practices self-attestation. We stress
that these are engineering-quality signals, not legal or compliance
certifications, and the platform's public documentation states this explicitly.

We report these facts as implementation context, not as an assessment of legal,
semantic, or model-analysis correctness. The per-attempt check-log schema was
designed so that the evaluation of Section~\ref{sec:eval} can be reconstructed
from operational data.

\section{Operational Evaluation of the Running Platform}
\label{sec:eval}

This section reports a reproducible, point-in-time operational measurement of
the running public platform. It deliberately answers narrower questions than a
semantic benchmark: what fraction of the configured inventory is exposed through
the public-evidence gate; which retrieval path is recorded for the current
inventory; and whether a seeded fixture instance exposes its placeholder records.
It does \emph{not} measure the correctness of a detected policy change, the
correctness of legal or model-generated interpretation, or latency across repeated
runs.

\subsection{Protocol and data artifact}

The production snapshot was collected at 2026-07-30 10:06:57 UTC from four
unauthenticated JSON endpoints of the running platform:
\url{/api/leaderboard}, \url{/api/source-suspensions},
\url{/api/changes?page=1&pageSize=50}, and \url{/api/companies}. The collection
capture record, endpoint digests, and interpretation notes are included with the
paper source. They access no
administrative API, credentials, subscriber information, private check logs, or
raw policy text. A local development instance with the seeded fixture inventory
was queried with the same endpoints as a control for the public-data gate.

The data artifact distinguishes a currently available record from a record that
is merely configured. It records the last retrieval path in the public
leaderboard, not a repeated-trial success probability. In particular, a
\texttt{PolicyChange} record visible through the public API is counted as a
public change-record inventory item; this evaluation does not independently
validate that record as a provider's legal amendment. The production snapshot
describes the deployed database, while the local control describes the bundled
fixture database after the Beta 21 migration.

\subsection{Results}

Table~\ref{tab:eval} reports the public production state. Of the 50 policy
records reported by the leaderboard, 45 (90.0\%) satisfied the public-evidence
gate and five (10.0\%) were suspended. The public companies endpoint returned 15
company records. The leaderboard reported 18 company rows, reflecting a broader
inclusion rule; because the release does not yet publish a versioned scope
manifest reconciling those denominators, we do not calculate a provider-level
coverage percentage. Four suspended documents belonged to Revolut and covered
its EU and UK privacy-policy and terms-of-use variants; the fifth was Amazon's EU
privacy notice. Their public representation exposed a minimal suspension state
but no policy text, score, timeline item, or analysis.

\begin{table}[htbp]
\centering
\small
\begin{tabularx}{\textwidth}{@{}l r r X@{}}
\toprule
\textbf{Public operational measure} & \textbf{Production} & \textbf{Seeded control} & \textbf{Interpretation} \\
\midrule
Policy records & 50 & 50 & Policy-row cardinality reported by each instance. \\
Public-evidence records & 45 (90.0\%) & 0 (0.0\%) & Records admitted to the public data gate. \\
Suspended records & 5 (10.0\%) & 50 of 50 & Minimal notices only; no policy text or analysis is exposed. \\
Public companies endpoint & 15 & 0 & Company records returned after public-data filtering; not a configured-scope denominator. \\
Public change-record inventory & 35 & 0 & Inventory of records that passed the public change gate, not a semantic-accuracy score. \\
\bottomrule
\end{tabularx}
\caption{Dated operational snapshot from the production public APIs and a local
seeded-fixture control. The two instances have the same configured inventory but
different evidence states.}
\label{tab:eval}
\end{table}
\FloatBarrier

The retrieval-path inventory contains 40 direct, five VPS-rendered, and five
archive-labeled policy rows. These are current inventory labels, not repeated
trial success rates. The public summary reports five renderer-backed and five
archive-backed rows but does not expose enough row-level evidence to infer that
every archive-labeled row is currently admissible for publication; the five
suspended records remain outside the public policy dataset.

The seeded-fixture control produced zero public-evidence records, zero public
change records, and zero publicly listed providers while returning 50 suspension
notices. This contrast is evidence that the tested public endpoints respect the
fixture-versus-evidence boundary under the supplied fixture state. It is not a
complete proof that every present or future output channel is free of leakage;
the code-level public-data gate and its test suite provide the complementary
engineering evidence described in Section~\ref{sec:impl}.

\subsection{What the snapshot does not establish}

The measurement has no independently annotated corpus of known policy revisions,
no expert legal coding, no repeated scan window, and no controlled baseline
system. It therefore cannot estimate precision or recall for change detection,
agreement with legal experts, robustness over time, or comparative advantage over
other monitoring systems. The 35 public change records are useful for inventory and interface
testing, but should not be read as 35 independently verified legal events. A
rigorous next study requires a versioned corpus of provider-published
texts, timestamped ground truth, blinded expert annotation, and a multi-week
repeated-scan protocol.

\section{Deployment and Operations}
\label{sec:deployment}

A design constraint throughout was operability by a single individual at consumer
cost. The web application, comprising the dashboard, the APIs, and the scanner,
runs on commodity shared hosting without root access. Headless rendering, which
requires a Chromium runtime, runs as an authenticated service of a few hundred
lines on a small virtual private server, reached over HTTPS and treated as
optional at every call site. The renderer refuses to start without its bearer
secret, applies the same address policy as the main application to every
subresource, blocks heavy subresources such as images and fonts for speed, uses a
fresh browser context per request to avoid cookie bleed, and bounds concurrency.
A separate operations agent on the same server performs backup, health probing,
update, and rollback. Scans are batchable, with bounded batch sizes and
per-company selection, to respect shared-hosting execution limits, and per-domain
politeness delays together with retry backoff bound the request rate to any
single provider. The production database is located outside the extracted
application directory so that redeploying the application cannot overwrite live
data, a safeguard added in direct response to the self-assessment of
Section~\ref{sec:redteam}. The resulting monthly operating cost is on the order of
a few tens of euros, which we consider part of the contribution: longitudinal
policy accountability does not require institutional infrastructure.

\section{Ethical Considerations}
\label{sec:ethics}

\paragraph{Purpose and proportionality.}
The system retrieves only public legal documents that providers are already
obligated to make available to their users. It does so at low frequency, on the
order of weekly scans of 50 documents, with politeness delays. It collects no
personal data and exercises no platform functionality beyond reading published
pages.

\paragraph{The fingerprinting tension.}
Reliable retrieval currently requires browser-realistic request fingerprints,
because several providers block identified crawlers from the very pages that bind
their users. That fact is itself a transparency finding. We document the tension
rather than hide it, and the roadmap includes an honest-first mode that first
attempts retrieval with an identifying agent and a contact URL and escalates to
browser-realistic fingerprints only when refused, together with the publication
of per-provider accessibility statistics.

\paragraph{Archives and load.}
Archive fallbacks impose load on public-interest infrastructure. The freshness
guard has the side effect of minimizing archive queries, and archive-derived data
is labeled with its capture time so that readers can judge its currency.

\paragraph{Publication integrity.}
The governance layer reflects an ethical stance. For a platform whose purpose is
trust, publishing nothing is strictly better than publishing plausible
fabrication. Every user-facing analysis carries its provenance and links to the
official document, and the analysis layer is presented to users as
AI-assisted rather than authoritative.

\section{Limitations and Future Work}
\label{sec:limits}

\paragraph{Evaluation.}
Section~\ref{sec:eval} reports a dated operational snapshot and a seeded-fixture
control. It does not report retrieval success rates \emph{over time},
false-change rates, or agreement between model-generated risk annotations and
expert legal coding. The check-log schema can support the first two measures; the
next study should pre-register a time window, preserve a provider-versioned
reference corpus, and include blinded expert review of a sampled set of detected
and non-detected changes.

\paragraph{Long-document analysis.}
Truncation before model analysis can miss clauses deep in long documents. We are
implementing a chunked map-reduce pipeline with an explicit coverage manifest
that records which spans of the document each conclusion derives from, enabling
audits for missed clauses and for conclusions unsupported by the text.

\paragraph{Raw evidence retention.}
The system stores normalized text. Retaining compressed raw HTML per snapshot,
with retention tiers, would strengthen the evidentiary chain and permit
retrospective re-extraction as the normalizer improves.

\paragraph{Extraction robustness.}
Accordion widgets, tabbed legal hubs, and multi-document pages remain extraction
hazards. Fragment-scoped extraction mitigates but does not eliminate them, and a
per-provider extraction test corpus with reference outputs is in progress.

\paragraph{Coverage and scale.}
Fifty documents across sixteen providers is a deliberate depth-over-breadth
choice. Scaling to hundreds of providers requires per-domain scheduling,
conditional requests using entity tags and modification times, and automated URL
rediscovery that follows legal-hub links and sitemaps when a configured URL dies.

\section{Conclusion}
\label{sec:conclusion}

\pw{} shows how continuous monitoring of the legal documents that govern consumer
technology can be organized around retrieval provenance and explicit withholding
states at civic cost. The mechanisms reported here, namely the freshness-guarded
retrieval cascade, the deterministic failure contract, the fixture-versus-evidence
boundary enforced at public egress, the evidence-gated re-baseline protocol, and
the adversarial questions that structure the self-assessment, are reusable design
patterns for systems that turn scraped web content into public records. The
operational snapshot is intentionally modest: it records what the running
platform exposed and withheld on a specific date. The source-reliability control
plane makes scan scope, acquisition reuse, failure causes, and remediation state
observable without converting them into public quality scores. Evidence Packets,
Collections, change-event polling, and receiver fixtures demonstrate how the
same provenance boundary can continue into reuse and integration. Semantic
accuracy, legal interpretation, adoption, and comparative performance remain
questions for a future expert-annotated, longitudinal study.

\subsection*{Availability}
The platform is publicly accessible at \url{https://policywatcher.online}.
The source code is released as open source and is subject to continuous
integration and third-party quality attestation as described in
Section~\ref{sec:impl}.

\begin{thebibliography}{99}

\bibitem{mcdonald2008cost}
A.~M. McDonald and L.~F. Cranor.
The cost of reading privacy policies.
\emph{I/S: A Journal of Law and Policy for the Information Society},
4(3):543-568, 2008.

\bibitem{obar2020biggest}
J.~A. Obar and A.~Oeldorf-Hirsch.
The biggest lie on the internet: ignoring the privacy policies and terms of
service policies of social networking services.
\emph{Information, Communication and Society}, 23(1):128-147, 2020.

\bibitem{fabian2017readability}
B.~Fabian, T.~Ermakova, and T.~Lentz.
Large-scale readability analysis of privacy policies.
In \emph{Proc. Int. Conf. on Web Intelligence (WI)}, 2017.

\bibitem{reidenberg2015disagreeable}
J.~R. Reidenberg, T.~Breaux, L.~F. Cranor, B.~French, et al.
Disagreeable privacy policies: mismatches between meaning and users'
understanding.
\emph{Berkeley Technology Law Journal}, 30(1):39-88, 2015.

\bibitem{wilson2016creation}
S.~Wilson, F.~Schaub, A.~A. Dara, F.~Liu, S.~Cherivirala, P.~G. Leon, et al.
The creation and analysis of a website privacy policy corpus.
In \emph{Proc. 54th Annual Meeting of the Association for Computational
Linguistics (ACL)}, 2016.

\bibitem{amos2021privacy}
R.~Amos, G.~Acar, E.~Lucherini, M.~Kshirsagar, A.~Narayanan, and J.~Mayer.
Privacy policies over time: curation and analysis of a million-document dataset.
In \emph{Proc. The Web Conference (WWW)}, 2021.

\bibitem{linden2020landscape}
T.~Linden, R.~Khandelwal, H.~Harkous, and K.~Fawaz.
The privacy policy landscape after the GDPR.
\emph{Proceedings on Privacy Enhancing Technologies (PoPETs)}, 2020(1):47-64,
2020.

\bibitem{zimmeck2014privee}
S.~Zimmeck and S.~M. Bellovin.
Privee: an architecture for automatically analyzing web privacy policies.
In \emph{Proc. 23rd USENIX Security Symposium}, 2014.

\bibitem{harkous2018polisis}
H.~Harkous, K.~Fawaz, R.~Lebret, F.~Schaub, K.~G. Shin, and K.~Aberer.
Polisis: automated analysis and presentation of privacy policies using deep
learning.
In \emph{Proc. 27th USENIX Security Symposium}, 2018.

\bibitem{srinath2021privacy}
M.~Srinath, S.~Wilson, and C.~L. Giles.
Privacy at scale: introducing the PrivaSeer corpus of web privacy policies.
In \emph{Proc. 59th Annual Meeting of the Association for Computational
Linguistics (ACL)}, 2021.

\bibitem{andow2019policylint}
B.~Andow, S.~Y. Mahmud, W.~Wang, J.~Whitaker, W.~Enck, B.~Reaves, K.~Singh, and
T.~Xie.
PolicyLint: investigating internal privacy policy contradictions on Google Play.
In \emph{Proc. 28th USENIX Security Symposium}, 2019.

\bibitem{cui2023poligraph}
H.~Cui, R.~Trimananda, A.~Markopoulou, and S.~Jordan.
PoliGraph: automated privacy policy analysis using knowledge graphs.
In \emph{Proc. 32nd USENIX Security Symposium}, 2023.

\bibitem{zimmeck2019maps}
S.~Zimmeck, P.~Story, D.~Smullen, A.~Ravichander, Z.~Wang, J.~Reidenberg,
N.~C. Russell, and N.~Sadeh.
MAPS: scaling privacy compliance analysis to a million apps.
\emph{Proceedings on Privacy Enhancing Technologies (PoPETs)}, 2019(3):66-86,
2019.

\bibitem{lippi2019claudette}
M.~Lippi, P.~Pa\l{}ka, G.~Contissa, F.~Lagioia, H.-W. Micklitz, G.~Sartor, and
P.~Torroni.
CLAUDETTE: an automated detector of potentially unfair clauses in online terms
of service.
\emph{Artificial Intelligence and Law}, 27(2):117-139, 2019.

\bibitem{utz2019consent}
C.~Utz, M.~Degeling, S.~Fahl, F.~Schaub, and T.~Holz.
(Un)informed consent: studying GDPR consent notices in the field.
In \emph{Proc. ACM SIGSAC Conference on Computer and Communications Security
(CCS)}, 2019.

\bibitem{nouwens2020dark}
M.~Nouwens, I.~Liccardi, M.~Veale, D.~Karger, and L.~Kagal.
Dark patterns after the GDPR: scraping consent pop-ups and demonstrating their
influence.
In \emph{Proc. CHI Conference on Human Factors in Computing Systems}, 2020.

\bibitem{matte2020cookie}
C.~Matte, N.~Bielova, and C.~Santos.
Do cookie banners respect my choice? Measuring legal compliance of banners from
IAB Europe's transparency and consent framework.
In \emph{Proc. IEEE Symposium on Security and Privacy (S\&P)}, 2020.

\bibitem{sanchezrola2019optout}
I.~Sanchez-Rola, M.~Dell'Amico, P.~Kotzias, D.~Balzarotti, L.~Bilge,
P.-A. Vervier, and I.~Santos.
Can I opt out yet? GDPR and the global illusion of cookie control.
In \emph{Proc. ACM Asia Conference on Computer and Communications Security
(AsiaCCS)}, 2019.

\bibitem{libert2018auditing}
T.~Libert.
An automated approach to auditing disclosure of third-party data collection in
website privacy policies.
In \emph{Proc. The Web Conference (WWW)}, 2018.

\bibitem{jueckstock2021realistic}
J.~Jueckstock, S.~Sarker, P.~Snyder, A.~Beggs, P.~Papadopoulos, M.~Varvello,
B.~Livshits, and A.~Kapravelos.
Towards realistic and reproducible web crawl measurement.
In \emph{Proc. The Web Conference (WWW)}, 2021.

\bibitem{iqbal2021fingerprinting}
U.~Iqbal, S.~Englehardt, and Z.~Shafiq.
Fingerprinting the fingerprinters: learning to detect browser fingerprinting
behaviors.
In \emph{Proc. IEEE Symposium on Security and Privacy (S\&P)}, 2021.

\bibitem{tosdr}
Terms of Service; Didn't Read.
\url{https://tosdr.org}. Accessed July 2026.

\bibitem{tosback}
Electronic Frontier Foundation.
TOSBack: the terms-of-service tracker.
\url{https://tosback.org}. Accessed July 2026.

\bibitem{wayback}
Internet Archive.
The Wayback Machine.
\url{https://web.archive.org}. Accessed July 2026.

\bibitem{commoncrawl}
Common Crawl Foundation.
The Common Crawl corpus.
\url{https://commoncrawl.org}. Accessed July 2026.

\bibitem{gdpr}
European Parliament and Council.
Regulation (EU) 2016/679 (General Data Protection Regulation), Articles 12 to 14.
\emph{Official Journal of the European Union}, L119, 2016.

\bibitem{dsa}
European Parliament and Council.
Regulation (EU) 2022/2065 (Digital Services Act), Article 14.
\emph{Official Journal of the European Union}, L277, 2022.

\bibitem{aiact}
European Parliament and Council.
Regulation (EU) 2024/1689 laying down harmonised rules on artificial
intelligence, Article 50.
\emph{Official Journal of the European Union}, L, 2024/1689, 2024.

\bibitem{ecarticle50}
European Commission.
Guidelines on the implementation of the transparency obligations for certain AI
systems under Article 50 of the AI Act.
Shaping Europe's Digital Future, 20 July 2026.
\url{https://digital-strategy.ec.europa.eu/en/library/guidelines-transparency-obligations-providers-and-deployers-ai-systems}.
Accessed 30 July 2026.

\bibitem{owasp2021}
Open Worldwide Application Security Project.
OWASP Top 10:2021, A10 Server-Side Request Forgery.
\url{https://owasp.org/Top10/}. Accessed July 2026.

\end{thebibliography}

\end{document}
